\section{Introducción}\label{sec:intro}

La evapotranspiración real ($\ETa$) es la componente del ciclo hidrológico
que devuelve a la atmósfera la mayor parte de la precipitación continental
y, en términos energéticos, el flujo de calor latente que cierra el balance
de energía de la superficie terrestre \parencite{fisher2017future}. En
cuencas agrícolas y forestales es, además, la medida física del uso
consuntivo del agua: la fracción que no retorna a cauces ni acuíferos. En
Chile central, la megasequía iniciada en 2010
\parencite{garreaud2017megadrought,garreaud2020megadrought} y la reforma del
Código de Aguas de 2022 \parencite{ley21435} han desplazado la discusión
pública desde el agua asignada hacia el agua efectivamente consumida. Un
mapa de $\ETa$ es, por ello, un insumo para balances hídricos de cuenca, la
fiscalización del uso consuntivo, la evaluación de programas de riego y la
valoración de servicios ecosistémicos de regulación hídrica. La utilidad de
ese insumo depende, sin embargo, de que su incertidumbre sea conocida.

Los modelos de balance de energía superficial alimentados por observaciones
de la Tierra estiman $\ETa$ como residuo del balance
$\Rn = G + H + LE$ \parencite{kalma2008review}. SEBAL
\parencite{bastiaanssen1998sebal1,bastiaanssen1998sebal2} y su derivado
METRIC \parencite{allen2007metric1,allen2007metric2} evitan la medición
directa del gradiente de temperatura aire--superficie mediante una
calibración interna: en dos píxeles ancla, uno frío y bien regado y otro
caliente y seco, se imponen condiciones de borde que fijan una relación
lineal entre ese gradiente y la temperatura superficial. METRIC ancla además
la calibración a la evapotranspiración de referencia de alfalfa ($\ETr$) y
extrapola al día mediante la fracción $\ETrF = ET_{inst}/\ETr$
\parencite{allen2007metric1,allen2011sebalmetric}. En la mayoría de las
aplicaciones operativas, el forzamiento meteorológico ---temperatura y
humedad del aire, viento y $\ETr$--- proviene de una única estación
considerada representativa de toda la escena.

Ese supuesto es débil en territorios heterogéneos. El forzamiento entra en
METRIC por al menos cinco vías: la radiación de onda larga incidente
(temperatura del aire), la transmisividad atmosférica (presión de vapor), el
viento a la altura de mezcla (velocidad del viento), la condición del ancla
fría ($\ETr$ instantánea) y la extrapolación diaria ($\ETr$ de 24~h). Parte
de los sesgos se absorbe en la calibración interna, de modo que su efecto
sobre $\ETa$ no es evidente \emph{a priori}
\parencite{long2011sebalsensitivity,long2013endmember}. La regresión
geográficamente ponderada (GWR) permite representar variables
meteorológicas como superficies cuya relación con predictores como la
elevación varía en el espacio \parencite{brunsdon1996gwr,fotheringham2002gwr},
pero su ventaja frente a interpoladores más simples no está garantizada:
depende de la densidad de la red, de la heterogeneidad del fenómeno y de la
estabilidad de los coeficientes locales
\parencite{wheeler2005multicollinearity,li2014interpolation}. Una mejor
interpolación de la temperatura del aire tampoco implica, por sí sola, una
mejor $\ETa$, porque la calibración interna de METRIC puede compensar o
amplificar las diferencias.

El software disponible no facilita contrastar esas preguntas. El paquete de
R \pkg{water} \parencite{olmedo2016water} implementa METRIC para Landsat con
una estación meteorológica; otras implementaciones son cerradas o están
acopladas a un sensor y a un forzamiento fijo. EarthFlux
\parencite{fuentesjaque2026earthflux} es un paquete de R de código abierto
que implementa la cadena METRIC como funciones independientes del sensor,
con forzamiento meteorológico intercambiable ---estación representativa,
regresión global, IDW o GWR--- bajo una interfaz común, validación cruzada
espacial y diagnósticos explícitos en lugar de correcciones silenciosas.

Este trabajo usa EarthFlux para responder la pregunta: \emph{¿en qué medida
la representación espacial de las condiciones meteorológicas mediante GWR
modifica la $\ETa$ estimada con METRIC y MODIS, en comparación con
configuraciones meteorológicas convencionales?} Se contrastan cuatro
hipótesis: (H1) los métodos de regionalización producen diferencias
espaciales identificables en las variables meteorológicas de METRIC; (H2)
esas diferencias se propagan de manera no uniforme hacia los componentes del
balance y la $\ETa$; (H3) la utilidad de GWR depende de la densidad de
estaciones, la heterogeneidad espacial y la configuración del modelo; y (H4)
una mayor complejidad metodológica no garantiza mayor exactitud. Se
distinguen expresamente cuatro niveles de evidencia que suelen confundirse:
la verificación de la implementación de las ecuaciones, las pruebas de
software, la validación de las variables meteorológicas regionalizadas y la
validación empírica de $\ETa$. Las contribuciones son: (i) una
implementación verificable y reproducible de METRIC con forzamiento espacial;
(ii) una evaluación fuera de muestra, con validación cruzada espacial
anidada, de la regionalización de las variables de METRIC en la región del
Maule; (iii) la cuantificación de cómo esas diferencias se propagan a
$\Rn$, $G$, $H$, $LE$ y $\ETa$ en una escena real, separando el efecto de
cada vía de entrada; y (iv) la comparación de ese efecto con la
incertidumbre de la calibración y de los insumos satelitales.
